\NativeSlot{44}{43}{%
% p010_s01

\ExamQuestion{7（续）}{5}%
\mbox{（\hspace{1.8em}）}，这个性质在修路时保证了路面宽度一致。\par

% p010_s02

\ExamQuestion{8.}{5}%
军军21时整准时睡觉，这时钟面上时针与分针所夹的角是\mbox{（\hspace{1.8em}）}度，我们也可以说时针与分针互相\mbox{（\hspace{1.8em}）}。\par
}
\NativeSlot{120}{43}{%
\ExamGroup{二、选择（将正确答案的序号填在括号里）。{\normalfont\kaishu\fontsize{10.8bp}{14bp}\selectfont （每题2分，共20分）}}

% p010_s03

\ExamQuestion{1.}{5}%
下列图形中，只有一组对边平行的是\mbox{（\hspace{1.8em}）}。\par
\ExamOptions{A. 长方形}{B. 平行四边形}{C. 梯形}{D. 六边形}
}
\NativeSlot{197}{121}{%
% p010_s04

\ExamQuestion{2.}{5}%
右图中\(C\)点到\(AB\)的距离是\mbox{（\hspace{1.8em}）}。\par
\ExamOptions{A. 3厘米}{B. 4厘米}{C. 5厘米}{D. 不确定}
\BookFigure{figures/p010_s04-01.tex}

% p010_s05

\ExamQuestion{3.}{5}%
下列说法正确的是\mbox{（\hspace{1.8em}）}。\par
\ExamOptions{A. 垂直的两条直线不一定相交}{B. 平行线无限延长也不会相交}{C. 两条线段平行，它们的长度一定相等}{D. 以上说法都对}
}
\NativeSlot{398}{36}{%
% p010_s06

\ExamQuestion{4.}{5}%
通过直线外一点可以画\mbox{（\hspace{1.8em}）}条直线与已知直线平行。\par
\ExamOptions{A. 1}{B. 2}{C. 3}{D. 无数}
}
\NativeSlot{447}{44}{%
% p010_s07

\ExamQuestion{5.}{5}%
直线行驶的汽车留下的车轮印迹\mbox{（\hspace{1.8em}）}。\par
\ExamOptions{A. 互相平行}{B. 互相垂直}{C. 相交}{D. 互相平行或互相垂直}
}
\NativeSlot{522}{36}{%
% p010_s08

\ExamQuestion{6.}{5}%
直线\(m\perp n\)，\allowbreak{}\(n\perp l\)，直线\(m\)和\(l\)的位置关系是\mbox{（\hspace{6em}）}。\par
\ExamOptions{A. 互相平行}{B. 互相垂直}{C. 相交}{D. 不确定}
}
\NativeSlot{571}{43}{%
% p010_s09

\ExamQuestion{7.}{5}%
下列物品中，两条边互相垂直的是\mbox{（\hspace{1.8em}）}。\par
\ExamOptions{A. 红领巾}{B. 三角尺}{C. 平行四边形装饰画}{D. 六边形杯垫}
}
\NativeSlot{646}{58}{%
% p010_s10

\ExamQuestion{8.}{5}%
下列关于过点\(A\)画已知直线\(l\)的垂线，步骤顺序正确的是\mbox{（\hspace{1.8em}）}。①把三角板一条直角边与直线\(l\)重合；②标出垂直符号；③沿着三角板另一条直角边画直线；④平移三角板，使直角顶点与点\(A\)重合。\par
\ExamOptions{A. ①②③④}{B. ②①③④}{C. ①④③②}{D. ③①②④}
}
\NativeSlot{744}{36}{%
% p010_s11

\ExamQuestion{9.}{5}%
在两条平行的步道之间修最短的人行道，该怎么画图纸？\mbox{（\hspace{1.8em}）}\par
\ExamOptions{A. 画一条斜线段}{B. 画一条垂线段}{C. 画一条任意线段}{D. 不确定}
}
